\chapter{Introdução}
\label{chap:introducao}

\orientacao{Diferente da introdução de uma macroentrega intermediária, a do relatório final abre no problema, porque será lida por quem não acompanhou o projeto. Cinco parágrafos, cada um com uma função:
\begin{enumerate}
    \item o problema, no domínio de aplicação, com a limitação das abordagens correntes;
    \item o desafio técnico específico e a lacuna que o projeto atacou;
    \item a hipótese ou a abordagem do projeto, em uma frase;
    \item os dados contratuais: período, duração, parceiros, arranjo de financiamento (inclusive o apoio da EMBRAPII) e o recorte de escopo, dizendo que o projeto entrega protótipo para validar métodos e algoritmos, e não produto. Essa frase custa uma linha e protege a unidade contra a expectativa de produto comercial;
    \item os pacotes de trabalho, em uma frase cada, a síntese dos resultados alcançados e a organização do documento.
\end{enumerate}}

\ph{Parágrafo 1: o problema no domínio de aplicação e a limitação das abordagens correntes.}

\ph{Parágrafo 2: o desafio técnico específico e a lacuna que o projeto atacou.}

\ph{Parágrafo 3: a hipótese ou a abordagem do projeto, em uma frase.}

O projeto \ph{NOME DO PROJETO} foi realizado pela \unidade{} com apoio da EMBRAPII, em parceria com \ph{parceiro tecnológico}, entre \ph{mm/aaaa} e \ph{mm/aaaa}. \ph{Recorte de escopo: o projeto entrega protótipo para validar métodos e algoritmos, e não produto.}

\ph{Parágrafo 5: os pacotes de trabalho, uma frase cada, e a síntese dos resultados.} O Capítulo~\ref{chap:objetivos} apresenta os objetivos; o Capítulo~\ref{chap:metodologia} descreve como o projeto foi conduzido; o Capítulo~\ref{chap:requisitos_final} consolida os requisitos; o Capítulo~\ref{chap:wp1} trata \ph{do pacote de trabalho 1}; e o Capítulo~\ref{chap:conclusao} reúne a síntese, as limitações e os trabalhos futuros.
